\begin{sectionbox}
    \subsection{Fourierpolynom}
    \begin{description}
        \item[Periode] betragsmäßig kleinste Zahl, wo sich die Funktion $\forall t$ wiederholt,
              z.B. $f(t)=\cos(k\omega t) + i\sin(k\omega t)$ wiederholt sich für $|k|\omega T_k=2\pi \Rightarrow T_k=\frac{2\pi}{|k|\omega}$
        \item[T-periodisch] Falls $f(t + T)=f(t) \quad \forall t\in\mathbb R$
        \item[Kreisfrequenz] $\omega = \frac{2\pi}{T}$ mit Periode $T$
    \end{description}

    Sei $\omega > 0, c_k \in \mathbb C, -n \le k \le n, n \in \mathbb N$. Dann heißt
    \begin{equation*}
        S_n(t)=\sum_{k=-n}^{n}c_ke^{ik\omega t}, t\in\mathbb R
    \end{equation*}
    Fourierpolynom bzw. trigonometrisches Polynom.

    \subsubsection{Eigenschaften}
    \begin{itemize}
        \item $T_1$ ist die größte auftretende Periode aller Summanden $\Rightarrow S_n$ ist $T_1$-periodisch
        \item $T_k= \frac{T_1}{|k|}$
        \item $\omega$ ist für alle Summanden gleich, $k$ variiert
    \end{itemize}
\end{sectionbox}

\subsection{Fourierreihe}
Jedes Fourierpolynom ist gleichzeitig eine Fourierreihe, wenn man alle nicht benötigten $c_k=0$ setzt.
\begin{emphbox}
    \begin{equation*}
        S_f(t)=\lim_{n\rightarrow\infty}S_n(t)=\sum_{k\in\mathbb Z}c_ke^{ik\omega t}=\sum_{k=-\infty}^{\infty}c_ke^{ik\omega t}
    \end{equation*}
\end{emphbox}

\begin{sectionbox}
    \subsubsection{Bestimmung der Fourierkoeffizienten}
    Sei eine $T$-periodische Funktion $f: \mathbb R \rightarrow \mathbb C$ gegeben und $k\in\mathbb Z$. Dann heißt
    \begin{equation*}
        c_k=\frac{1}{T}\int_a^{a+T} f(t)e^{-ik\omega t}\diff t, \quad a \text{ bel.}
    \end{equation*}
    der $k$-te Fourierkoeffizient von $f$. \\
    \textbf{Fourierspektrum} Familie $\{c_k\}_{k\in\mathbb Z}$ aller Fourierkoeffizienten. Beschreibt Amplituden und Phasenverschiebungen der Teilschwingungen in $f$.

    \subsubsection{Konvergenz der Fourierreihe}
    \begin{sectionbox}
        $f:\mathbb R \rightarrow \mathbb C$ ist stückweise stetig, falls $\forall t \in \mathbb R$ die rechts- und linksseitigen Grenzwerte existieren und $f$ auf jedem beschränkten Interval nur in endlich vielen Punkten unstetig ist (Sprungstellen)
    \end{sectionbox}

    \begin{sectionbox}
        $f$ stückweise stetig ist stückweise stetig differenzierbar, falls $f$ auf jedem beschränkten Intervall nur in endlich vielen Punkten nicht stetig differenzierbar ist und die rechts- und linksseitigen Grenzwerte von $f'$ auch existieren
    \end{sectionbox}

    Sei $f$ stückweise stetig differenzierbar:
    \begin{description}
        \item[Punktweise Konvergenz] Fourierreihe konvergiert $\forall t$ gegen $\frac{f(t^-) + f(t^+)}{2}$
        \item[Gleichmäßige Konvergenz] Auf jedem abgeschlossenen Intervall ohne Sprungstellen konvergiert die Fourierreihe gegen $f$
    \end{description}
\end{sectionbox}

\begin{sectionbox}
    \subsubsection{Umrechnung zu sin-/cos-Darstellung}
    \begin{emphbox}
        \begin{equation*}
            S_f(t)=\frac{a_0}{2}+\lim_{n\rightarrow\infty}\sum_{k=1}^{n}(a_k\cos(k\omega t) + b_k\sin(k\omega t))
        \end{equation*}
    \end{emphbox}
    mit $k\in\mathbb N$:
    \begin{equation*}
        a_0=2c_0 \quad a_k=c_k+c_{-k} \quad b_k=i(c_k-c_{-k})
    \end{equation*}
    bzw.
    \begin{align*}
        a_k=\frac{2}{T}\int_0^Tf(t)\cos(k\omega t)\diff t \quad k \in \mathbb N_0 \\
        b_k=\frac{2}{T}\int_0^Tf(t)\sin(k\omega t)\diff t \quad k \in \mathbb N
    \end{align*}

    Falls $f$ reel: $c_{-k} = \overline{c_k}$,
    \begin{equation*}
        a_0=2c_0 \quad a_k=2\Re{c_k} \quad b_k=-2\Im{c_k}
    \end{equation*}
\end{sectionbox}

\begin{sectionbox}
    Bei $T$-periodischem $f$: \\
    $f$ gerade $\Leftrightarrow f(\frac{T}{2}+t) = f(\frac{T}{2}-t)$ \\
    $f$ ungerade $\Leftrightarrow f(\frac{T}{2}+t) = -f(\frac{T}{2}-t)$

    Seien $f,g:\mathbb R \rightarrow\mathbb C \; T$-periodisch. Sei $\{c_k\}_{k\in\mathbb Z}$ das Fourierspektrum von $f$ und $\{d_k\}_{k\in\mathbb Z}$ von $g$
    \subsubsection{Eigenschaften}
    \begin{itemize}
        \item $f$ gerade $\Rightarrow c_k=c_{-k}$ bzw. $b_k=0 \; \forall k \in\mathbb N$. Fourierreihe besteht nur aus $\cos$-Termen
        \item $f$ ungerade $\Rightarrow c_k=-c_{-k}$ bzw. $a_k=0 \; \forall k \in\mathbb N_0$. Fourierreihe besteht nur aus $\sin$-Termen
        \item Das Fourierspektrum von $\alpha f(t) + \beta g(t)$ ist $\{\alpha c_k + \beta d_k\}_{k\in\mathbb Z}$
        \item Das Fourierspektrum von $\overline{f(t)}$ ist $\{\overline{c_{-k}}\}_{k\in\mathbb Z}$
        \item Das Fourierspektrum von $f(-t)$ ist $\{c_{-k}\}_{k\in\mathbb Z}$
        \item Das Fourierspektrum von $f(t + a)$ (Verschiebung im Zeitbereich) ist $\{e^{ik\omega a}c_k\}_{k\in\mathbb Z}$
        \item Das Fourierspektrum von $e^{in\omega t}f(t) \quad n\in\mathbb Z$ (Amplitudenmodulation) ist $\{c_{k-n}\}_{k\in\mathbb Z}$
        \item Das Fourierspektrum von $f(\gamma t)$ (Streckung) ist $\{c_k\}_{k\in\mathbb Z}$
        \item Sei $f$ stetig und $f'$ stückweise stetig: Das Fourierspektrum von $f'(t)$ ist $\{ik\omega c_k\}_{k\in\mathbb Z}$
        \item Sei f stückweise stetig und $c_0=0$: Das Fourierspektrum von $F(t)=\int_0^t f(s)\diff s$ ist
              \begin{equation*}
                  s_k=\begin{cases}
                      -\frac{1}{T}\int_0^Ttf(t)\diff t & k=0    \\
                      \frac{c_k}{ik\omega}             & k\ne 0
                  \end{cases}
              \end{equation*}
    \end{itemize}
\end{sectionbox}

\begin{sectionbox}
    \subsubsection{Orthogonalreihenrelationen}
    Für $k_1, k_2 \in \mathbb Z$ gilt:
    \begin{equation*}
        \frac{1}{T}\int_0^Te^{ik_1\omega t}e^{-ik_2\omega t}\diff t = \delta_{k_1,k_2}=\begin{cases}
            1, & k_1=k_2    \\
            0, & k_1\ne k_2
        \end{cases}
    \end{equation*}
    $\Rightarrow f_k(t)=e^{ik\omega t}$ hat Fourierspektrum $c_{k'}=\delta_{k,k'} \; \forall k' \in \mathbb Z$
\end{sectionbox}